Ihre Arbeit wird nach zehn gleichgewichteten, in vier Gruppen unterteilten Kriterien bewertet.

I.	Bearbeitung 
1.	Einarbeitung, Verständnis und Überblick unter Berücksichtigung des Schwierigkeitsgrades
2.	Einfallsreichtum, Ideen, Initiative; Selbständigkeit, Entscheidungsfreudigkeit
3.	systematisches Vorgehen, wissenschaftliche Eigenleistung
4.	Organisationstalent, Zeiteinteilung; Ausdauer, Fleiß
II.	Ergebnis (Theorie, Software, Hardware)
	5.	Inwieweit wurde das mögliche Ziel unter Berücksichtigung der
	Anforderungen erreicht?
	6.	Verwendbarkeit der Ergebnisse: Tragfähigkeit der Theorie und Methodik, Funktionstüchtig-keit der Hardware und der Software
III.	Schriftliche Ausarbeitung (Dokumentation)
	7.	Klarheit von Gliederung und Aufbau, Vollständigkeit und wiss. Korrektheit, Bewertung der Ergebnisse
	8.	Stil, Ausdruck; Prägnanz der Abbildungen und Diagramme, 
grafische Gestaltung
IV.	Abschlusspräsentation
	9.	Didaktischer Aufbau, wiss. Inhalt, kritische Bewertung;
	Diskussionsverhalten
10.	Vortragsstil, Zeitdisziplin; Vortragsmaterialien: Folien, Video, 
PowerPoint, etc.)

Jedes der 10 Kriterien wird mit einer Note zwischen 1 und 5 bewertet. Beachten Sie also, dass die geordne-te, exakte Darstellung Ihrer Arbeit und der Ergebnisse sowohl in der schriftlichen Ausarbeitung als auch im Abschlussvortrag zusammen mit 40% zur abschließenden Bewertung beitragen.


Bewertung des Abschlusskolloquiums
In der folgenden Checkliste sind noch einmal die Punkte aufgeführt, die bei der Bewertung Ihres Vortrags
berücksichtigt werden:
 Wie treten Sie in Ihrem Abschlusskolloquium auf?
 Wie lang ist Ihr Vortrag? Halten Sie sich an die Vorgabe von 20 Minuten?
 Wie ist der Vortrag strukturiert?
 Wie sind Ihre Folien gestaltet?
 Verwenden Sie ausreichend viele Bilder?
 Wie aussagekräftig sind die Bilder?
 Wie gehen Sie mit Fragen um?
 Wie gut können Sie die gestellten Fragen beantworten?
Der letzte Punkt ist besonders wichtig, weil Sie über ihn

\subsection{ba}
Die 10 Kriterien sind mit den ganzzahligen Noten 1 bis 5 zu bewerten. Erklärungen zu den einzelnen Noten bitte in Form eines kurzen Gutachtens ergänzen. (maximal eine DIN A4 Seite)
(„1“ sehr gut; „2“ gut; „3“ befriedigend; „4“ ausreichend; „5“ nicht ausreichend)
I. Bearbeitungsphase
Lfd.Nr.
Kriterien
Note
1
Einarbeitung, Verständnis und Überblick unter Berücksichtigung des Schwierigkeitsgrades
2
Einfallsreichtum, Ideen, Initiative, Selbständigkeit, Entscheidungsfreudigkeit
3
systematisches Vorgehen,
wissenschaftliche Eigenleistung
4
Organisationstalent, Zeiteinteilung;
Ausdauer, Fleiß
II. Ergebnis (Theorie, Software, Hardware)
5
Inwieweit wurde das mögliche Ziel unter Berücksichtigung der Anforderungen erreicht
6
Verwendbarkeit der Ergebnisse: Tragfähigkeit der Theorie und Methodik und/oder Funktionstüchtigkeit der Hardware und der Software
III. Schriftliche Ausarbeitung (Dokumentation)
7
Klarheit von Gliederung und Aufbau
8
Vollständigkeit und wissenschaftliche Korrektheit, Bewertung der Ergebnisse
9
Stil, Ausdruck
10
Prägnanz der Abbildungen und Diagramme, grafische Gestaltung

\begin{itemize}
    \item didaktischer Aufbau des Vortrags
    \item Vortrag: fachlicher Inhalt, kritische Bewertung der Ergebnisse
    \item Vortragsstil, Zeitdisziplin
    \item Qualität der Vortagsmaterialien: Folien, Video, Demonstratoren etc., wissenschaftliche Korrektheit
II. 
    \item Diskussionsverhalten bezüglich des Vortrags
    \item Diskussionsverhalten bezüglich der schriftlichen Ausarbeitung
    \item Fragen und Diskussion: fachliche Kompetenz (Ingenieurstechnische Grundlagen)
    \item Fragen und Diskussion: fachliche Kompetenz (Elektrotechnik)
    \item Fragen und Diskussion: fachliche Kompetenz (themenspezifisch)
    \item Fragen und Diskussion: methodische Kompetenz
\end{itemize}


Die praktische bzw. theoretische Arbeit nimmt den größten Teil der Zeit in Anspruch, die Sie für die Abschlussarbeit
benötigen. Demzufolge wird dieser Teil auch mit einem hohen Gewicht bewertet (siehe
oben). Hierbei werden die folgenden Punkte berücksichtigt:
\begin{itemize}
    \item Wie gut und wie schnell haben Sie sich in das gestellte Aufgabengebiet eingearbeitet?
\item ie groß war der Aufwand für die Einarbeitung?
\item ie gut wurde die Arbeit von Ihnen insgesamt fachlich bearbeitet?
\item Wie umfangreich ist das Wissen zu dem zu bearbeiteten Themenkomplex?
\item Wie kreativ waren Sie bei der Lösung auftretender Probleme?
\item Mit wie viel Energie haben Sie das Thema bearbeitet (Fleiß)?
\item Wie gut ist eine Lösung ausgearbeitet und wie elegant ist sie? (Hierzu gehört z.B. auch, dass ein
Programm ausreichend kommentiert sein muss, dass es strukturiert ist, dass Redundanzen vermieden
werden usw.)
\item Welche Fehler oder Defizite weist das Projekt auf?
\item Wurden alle Punkte bearbeitet, die in der Aufgabenstellung definiert waren?
Bei all diesen Bewertungskriterien spielt indirekt immer auch die Komplexität der Aufgabenstellung eine
Rolle, da ein einfaches Thema leichter zu lösen ist, als ein komplexes Thema.
\end{itemize}
